\section*{第 \ref{ch:b3:colloids} 章——界面、表面活性剂与胶体}

\begin{solution}{exo:b3:colloids:1}
$r = \qty{0.50}{\micro m}$：$\Delta p = 2 \times 0.0720/\num{5.0e-7} = \qty{2.9e5}{Pa}$；内部压强为 $1.0 + 2.9 = \qty{3.9}{bar}$。
\end{solution}

\begin{solution}{exo:b3:colloids:2}
$\ln(p/p^*) = 2 \times 0.0720 \times \num{1.807e-5}/(\num{1.0e-8} \times 8.314 \times 298.15) = 0.105$; $p/p^* = 1.11$.
\end{solution}

\begin{solution}{exo:b3:colloids:3}
$\cos\theta = (40 - 20)/72 = 0.278$，$\theta = \ang{74}$：小于 $90^\circ$，液体
部分\omterm{def:b3:colloids:wetting}{润湿}固体。
\end{solution}

\begin{solution}{exo:b3:colloids:4}
NaCl：$I = \qty{10}{mol.m^{-3}}$，$\kappa^{-1} = 0.96\sqrt{100/10} = \qty{3.0}{nm}$。\ce{MgSO4}：$I = \qty{40}{mol.m^{-3}}$，\qty{1.5}{nm}。
\end{solution}

\begin{solution}{exo:b3:colloids:5}
$\Gamma = \num{12.6e-3}/(2 \times 8.314 \times 298.15) = \qty{2.54e-6}{mol.m^{-2}}$；面积为 $1/(\Gamma N_A) =
\qty{0.65}{nm^2}$ 每个分子。
\end{solution}

\begin{solution}{exo:b3:colloids:6}
低于\omterm{def:b3:colloids:micelle}{临界胶束浓度}时，斜率为 $-12.9$ \unit{mN/m}（每单位 $\log c$，自 $-5.0$ 至 $-4.0$）；
平台值为 33.4。直线在 $\log c = -4.0 + (39.1 - 33.4)/12.9 = -3.56$ 处与平台相交，
即\omterm{def:b3:colloids:micelle}{临界胶束浓度} $\approx \qty{2.8e-4}{mol/L}$。
\end{solution}

\begin{solution}{exo:b3:colloids:7}
SDS：$P = 0.350/(0.60 \times 1.67) = 0.35$，位于球形与短圆柱的分界处，
\omterm{def:b3:colloids:micelle}{临界胶束浓度}附近形成球形\omterm{def:b3:colloids:micelle}{胶束}。脂质：$P = 0.700/(0.65 \times 1.67) = 0.64$：
形成双分子层，因而为囊泡。
\end{solution}

\begin{solution}{exo:b3:colloids:8}
\ce{CaCl2}：$60/64 = \qty{0.94}{mM}$（以 \ce{Ca^{2+}} 计）；\ce{AlCl3}：$60/729 = \qty{0.082}{mM}$（以 \ce{Al^{3+}} 计）。\ce{Na3PO4}
通过其中的 \ce{Na+} \omterm{def:b3:colloids:stability}{聚沉}（盐浓度约 \qty{20}{mM}）；磷酸根为
共离子，甚至可能被\omterm{def:b3:surfaces-catalysis:adsorption}{吸附}并使负电荷增加。
\end{solution}

\begin{solution}{exo:b3:colloids:9}
亲水部分为 $6 \times 44.0 + 17.0 = \qty{281.0}{g/mol}$，占 \qty{450.0}{g/mol} 的比例（取教材原子量）：
\omterm{def:b3:colloids:hlb}{亲水亲油平衡值} HLB $= 20 \times 281.0/450.0 = 12.5$，为水包油型乳化剂。
\end{solution}

\begin{solution}{exo:b3:colloids:10}
小液滴中的油在连续相中的溶解度更大（按因子
$\exp(2\gamma V_{\mathrm m}/rRT)$，\qty{50}{nm} 液滴的指数项比 \qty{500}{nm} 液滴大十倍）：它
扩散到大液滴上，大液滴长大而小液滴消失。这
种\omterm{def:b3:colloids:ripening}{奥斯特瓦尔德熟化}使\omterm{def:b3:colloids:colloid}{乳液}变粗，即使两个液滴从未合并。
\end{solution}

\begin{solution}{exo:b3:colloids:11}
弯月面为半径 $r/\cos\theta$ 的球冠：紧邻其下方的液体
压强低 $2\gamma\cos\theta/r$，液柱上升直至 $\rho gh$ 与之抵消。对于水，
$h = 2 \times 0.072/(997 \times 9.81 \times \num{1.0e-4}) = \qty{0.15}{m}$。
\end{solution}

\begin{solution}{exo:b3:colloids:12}
在 \qty{100}{mM} 下\omterm{def:b3:electrode-kinetics:double-layer}{德拜长度}小于 \qty{1}{nm}，排斥在吸引之前消失，
即已无势垒（见 \qty{50}{mM} 曲线之外的部分）。游离聚合物不在
表面起作用；接枝聚合物层无论盐浓度如何，在接触时都会排斥，
这是链重叠时所损失的熵。
\end{solution}

\begin{solution}{pb:b3:colloids:1}
\textbf{1.} 晶格中的置换（\ce{Al^{3+}} 替代 \ce{Si^{4+}}，\ce{Mg^{2+}} 替代 \ce{Al^{3+}}）留下
永久负电荷，由可交换阳离子所抵消。
\textbf{2.} $I = \frac12(1.0 + 1.0) = \qty{1.0}{mM}$.
\textbf{3.} $\kappa^{-1} = \qty{9.6}{nm}$.
\textbf{4.} $a\kappa = 10$：\omterm{def:b3:electrode-kinetics:double-layer}{双电层}相对于颗粒较薄，这是
德亚金近似的条件。
\textbf{5.} $v = 2 \times 1650 \times 9.81 \times (\num{1.0e-7})^2/(9 \times \num{0.89e-3}) = \qty{4.0e-8}{m/s}$，约合每天 \qty{3.5}{mm}。
\textbf{6.} 两者的\omterm{def:b3:electrode-kinetics:double-layer}{扩散层}先发生排斥，然后范德瓦耳斯吸引才起主导作用。
\textbf{7.} $V = 2\pi\varepsilon_r\varepsilon_0a\psi^2\ln(1 + \eu^{-\kappa h}) - Aa/12h$.
\textbf{8.} 约为 $39\,kT$。
\textbf{9.} 碰撞中所占比例的量级为 $\eu^{-39} \approx 10^{-17}$：实际上几乎从不发生。
\textbf{10.} 该势垒消失：每次碰撞都导致接触。
\textbf{11.} \omterm{def:b3:electrode-kinetics:double-layer}{德拜长度}越短，排斥作用局限于小间隙，
在这些间隙中，随 $1/h$ 增长的吸引作用占主导。
\textbf{12.} \omterm{def:b3:colloids:stability}{临界聚沉浓度}随抗衡离子电荷按 $z^{-6}$ 变化。
\textbf{13.} $50/64 = \qty{0.78}{mM}$.
\textbf{14.} $50/729 = \qty{0.069}{mM}$.
\textbf{15.} 阳离子是带负电颗粒的抗衡离子：它们
在\omterm{def:b3:electrode-kinetics:double-layer}{双电层}中聚集并起屏蔽作用。
\textbf{16.} 硫酸根为共离子，被颗粒排斥。
\textbf{17.} 每立方米含 \qty{0.069}{mol} 的 \ce{Al^{3+}}。
\textbf{18.} \qty{0.034}{mol} 的 $\ce{Al2(SO4)3}$。
\textbf{19.} $0.0343 \times 342.3 = \qty{11.7}{g}$.
\textbf{20.} \qty{11.7}{kg} 每小时。
\textbf{21.} $\ce{Al^3+ + 3 HCO3- -> Al(OH)3 + 3 CO2}$.
\textbf{22.} $3 \times 0.069 = \qty{0.21}{mM}$，占 \qty{1.0}{mM} 的五分之一。剩余的
碳酸氢根与溶解的 \ce{CO2} 对水起缓冲作用：pH
仅适度下降。
\textbf{23.} 高电荷的铝物种发生\omterm{def:b3:surfaces-catalysis:adsorption}{吸附}，使颗粒电荷反转，
颗粒之间重新相互排斥。
\textbf{24.} \textbf{每立方米约需硫酸铝 \qty{12}{g}}（\qty{11.7}{g}），
再加上实际水解所需的额外用量。
\end{solution}
